\section{Construcción de la planta (puntos extra)}
\label{sec:p6}
\enunciado{Construya la planta del proyecto (solo no tiene controlador pero si el resto de
componentes) (40\,\% extra). Adjunte fotos y video de la construcción final.}

\puntosextra{Se construyó la planta completa y se muestra funcionando en lazo abierto en el
video \videoPlanta. Solo falta cerrar el lazo con el PID, que corresponde al Proyecto Grupal~2.}

\subsection{Componentes y conexiones}
La Tabla~\ref{tab:planta} resume la planta construida, que implementa los bloques de la
Fig.~\ref{fig:sistema}.

\begin{table}[h]
\centering\small
\caption{Componentes y conexiones de la planta construida.}
\label{tab:planta}
\begin{tabularx}{\linewidth}{@{}>{\raggedright\arraybackslash}p{3.5cm}>{\raggedright\arraybackslash}X>{\raggedright\arraybackslash}p{3.3cm}@{}}
\toprule
Componente & Función & Conexión \\
\midrule
Disco de MDF & Rueda frenada; 10 marcas negras equidistantes sobre fondo blanco & Eje del motor \\
Motor CD & Lleva el disco a la velocidad inicial & Salidas OUT1/OUT2 del L298N \\
Puente H L298N & Etapa de potencia del motor & ENA 25, IN1 26, IN2 33 \\
Fuente del motor & Alimenta el motor a través del puente & Bornes +12V y GND del L298N \\
Servo MG90S + goma & Freno: 120° liberado, 180° frenado & Señal en GPIO18 (PWM 50~Hz) \\
Sensor TCRT5000 & Un pulso por marca, 10 por vuelta & Salida en GPIO27, a 3,3~V \\
ESP32 (ESP-IDF 5.5) & Mide, registra y manda freno y motor & USB a la computadora \\
\bottomrule
\end{tabularx}
\end{table}

\subsection{Funcionamiento del sistema}
\begin{itemize}
  \item \textbf{Motor CD y puente H.} Un pin del ESP32 entrega 3,3~V y pocos miliamperios,
  insuficientes para el motor, así que el L298N conmuta la corriente de una \textbf{fuente externa}
  conectada a sus bornes de potencia. El ESP32 solo envía señales lógicas: IN1/IN2 fijan el sentido
  y ENA, con PWM de 1~kHz, fija el porcentaje de voltaje aplicado. Para frenar, ENA pasa a bajo y el
  motor queda en \emph{marcha libre}. Así el motor no frena el disco y la única acción de frenado es la
  de la goma, que es la que se quiere caracterizar.
  \item \textbf{Freno.} El MG90S es un servo de posición: recibe un pulso cada 20~ms cuyo ancho
  (500 a 2500~µs) fija el ángulo. A 120° la goma está separada del disco y a 180° lo presiona. El
  freno es unilateral, como se discutió en la pregunta~5: solo quita velocidad.
  \item \textbf{Sensor óptico.} El TCRT5000 detecta por reflexión el paso de cada marca negra. El
  ESP32 fecha un solo flanco por marca con una interrupción y calcula la velocidad por tiempo entre
  marcas (pregunta~2). Un filtro descarta los flancos separados por menos de 500~µs, que son rebotes
  del comparador del módulo y no marcas reales.
  \item \textbf{ESP32 y computadora.} El firmware del banco (\texttt{firmware/} del repositorio)
  recibe comandos por la consola serie desde la computadora: \texttt{FRENO} mueve el servo,
  \texttt{MOTOR} impulsa el disco y \texttt{ENSAYO} ejecuta el ensayo de la pregunta~7. Durante una
  captura envía una fila \texttt{tiempo\_ms, freno\_deg, velocidad\_rpm} cada 20~ms.
  \item \textbf{Alimentación y tierra común.} El motor y el servo consumen picos de corriente que
  harían caer el voltaje del ESP32. Por eso cada potencia va por su lado y todas las tierras se unen,
  para que las señales lógicas tengan la misma referencia.
\end{itemize}

\subsection{Video en lazo abierto}
El video \videoPlanta{} muestra la planta operando sin controlador: el motor acelera el disco, se
corta, y el servo aplica la goma con un ángulo fijo, mientras el ESP32 mide la velocidad.

\IfFileExists{figuras/planta.jpg}{%
\begin{figure}[h]
\centering
\includegraphics[width=0.8\linewidth]{planta.jpg}
\caption{Planta construida.}
\label{fig:planta}
\end{figure}}{}
